\chapter{Nonlinear Diophantine Equations}

\begin{de}
Let $d\in\N$ be nonsquare.  The equations
\[
 x^2-dy^2=1
 \qquad\text{and}\qquad
 x^2-dy^2=-1
\]
are the positive and negative Pell equations.
\end{de}

\begin{lem}
Let $x_0,y_0\in\N$, let $\delta>0$ with $\sqrt\delta\notin\Q$, and suppose
\[
 x_0^2-\delta y_0^2=\alpha,
 \qquad
 0\le\alpha<\sqrt\delta.
\]
Then
\[
 \abs*{\frac{x_0}{y_0}-\sqrt\delta}<\frac1{2y_0^2}.
\]
\end{lem}
\begin{proof}
The case $\alpha=0$ would make $\sqrt\delta=x_0/y_0$ rational, so in fact $\alpha>0$.  Hence $x_0^2>\delta y_0^2$, $x_0>\sqrt\delta y_0$, and
\[
 \frac{x_0}{y_0}-\sqrt\delta
 =\frac{\alpha}{y_0(x_0+\sqrt\delta y_0)}
 <\frac{\sqrt\delta}{2\sqrt\delta y_0^2}.
\]
\end{proof}

\begin{thm}
Let $d$ be a positive nonsquare integer and let $n\in\Z$ satisfy $\abs n<\sqrt d$.  If $(x_0,y_0)\in\N^2$ satisfies
\[
 x_0^2-dy_0^2=n,
\]
then $x_0/y_0$ is a convergent of $\sqrt d$.
\end{thm}
\begin{proof}
If $n>0$, Lemma 7.1 and Legendre's criterion apply directly.  If $n<0$, rewrite the equation as
\[
 y_0^2-\frac1d x_0^2=-\frac nd,
 \qquad
 0<-\frac nd<\frac1{\sqrt d}.
\]
After reducing the fraction if necessary, Lemma 7.1 and Legendre's criterion show that $y_0/x_0$ is a convergent of $1/\sqrt d$.  The continued fraction of $1/x$ is obtained by adding a leading zero, so its nonzero convergents are the reciprocals of the convergents of $x$.  Hence $x_0/y_0$ is a convergent of $\sqrt d$.  The case $n=0$ is impossible because $d$ is nonsquare.
\end{proof}

For the rest of the Pell discussion, put
\[
 x_0=\floor{\sqrt d}+\sqrt d
\]
and let its purely periodic continued fraction have least period $r$.  Let
$(P_i,Q_i)$ be the complete-quotient data of Chapter 6, with $Q_0=1$.

\begin{thm}
\begin{enumerate}[label=(\alph*)]
\item $Q_i=1$ if and only if $i\equiv0\pmod r$;
\item $Q_i\ne-1$ for every $i$.
\end{enumerate}
\end{thm}
\begin{proof}
The number $x_0$ is reduced, so its complete quotients are reduced and repeat with least period $r$.  Hence $i\equiv0\pmod r$ implies $Q_i=1$.  Conversely, if $Q_i=1$, then
$x_i=P_i+\sqrt d$.  Reducedness gives
$-1<P_i-\sqrt d<0$, so $P_i=\floor{\sqrt d}$ and $x_i=x_0$.  Minimality of the period gives $r\mid i$.

If $Q_i=-1$, then
$x_i=-P_i-\sqrt d>1$, whereas its conjugate $x_i^*=-P_i+\sqrt d$ is positive, contradicting reducedness.
\end{proof}

\begin{thm}
Let $p_k/q_k$ be the convergents of $\sqrt d$ and let $r$ be its least period.
\begin{enumerate}[label=(\alph*)]
\item If $r$ is even, the negative Pell equation has no positive solution, and the positive solutions are
\[
 (p_{tr-1},q_{tr-1}),\qquad t=1,2,3,\ldots.
\]
\item If $r$ is odd, the negative solutions are the same pairs with $t$ odd, and the positive solutions are those with $t$ even.
\end{enumerate}
\end{thm}
\begin{proof}
By Theorem 7.2, a solution of $x^2-dy^2=\pm1$ must be a convergent.  Theorem 6.20 gives
\[
 p_k^2-dq_k^2=(-1)^{k+1}Q_{k+1}.
\]
Theorem 7.3 shows that this equals $\pm1$ exactly when $k+1=tr$.  The sign is then $(-1)^{tr}$, which gives the two cases.
\end{proof}

\begin{de}
Among the positive solutions to $x^2-dy^2=1$, the one with smallest $x$ is the fundamental solution.
\end{de}

\begin{thm}
Let $(x_1,y_1)$ be the fundamental solution of $x^2-dy^2=1$.  All positive solutions are given by
\[
 x_n+y_n\sqrt d=(x_1+y_1\sqrt d)^n,
 \qquad n\in\N.
\]
\end{thm}
\begin{proof}
The displayed powers have integer coefficients and norm $1$, so they are solutions.  Conversely, let
$\eta=s+t\sqrt d>1$ be any positive solution and put
$\varepsilon=x_1+y_1\sqrt d$.  Choose $n\ge0$ such that
\[
 \varepsilon^n\le\eta<\varepsilon^{n+1}.
\]
Then
\[
 \delta=\eta\varepsilon^{-n}=A+B\sqrt d
\]
has integer coefficients, norm $1$, and $1\le\delta<\varepsilon$.  Its conjugate is $\delta^{-1}>0$, so
$A=(\delta+\delta^{-1})/2>0$ and $B\ge0$.  If $B>0$, then $(A,B)$ is a positive Pell solution.  Since the function $z+z^{-1}$ is increasing for $z\ge1$,
\[
 A=\frac{\delta+\delta^{-1}}2
 <\frac{\varepsilon+\varepsilon^{-1}}2=x_1,
\]
contradicting the definition of the fundamental solution.  Thus $B=0$, whence $\delta=1$ and $\eta=\varepsilon^n$.
\end{proof}

\begin{app}
\begin{enumerate}[label=(\alph*)]
\item There are infinitely many triples of consecutive integers, each of which is a sum of two integer squares.  For every positive solution of $x^2-2y^2=1$,
\[
 x^2-1=y^2+y^2,
 \qquad
 x^2=x^2+0^2,
 \qquad
 x^2+1=x^2+1^2.
\]
\item If $x^2-dy^2=n$ has a nonzero integer solution $(s,t)$ and the positive Pell equation for $d$ has solutions $(x_m,y_m)$, then
\[
 (s+t\sqrt d)(x_m+y_m\sqrt d)=S_m+T_m\sqrt d
\]
gives infinitely many solutions $S_m^2-dT_m^2=n$.
\item There are infinitely many positive integers $n$ for which $n^2+(n+1)^2$ is a square.  Indeed,
\[
 n^2+(n+1)^2=m^2
 \quad\Longleftrightarrow\quad
 (2n+1)^2-2m^2=-1,
\]
and the negative Pell equation for $d=2$ has infinitely many solutions.  Geometrically, these are infinitely many integer right triangles with consecutive legs.
\end{enumerate}
\end{app}

\section{Pythagorean triples}

\begin{de}
A triple $(x,y,z)\in\N^3$ is Pythagorean if $x^2+y^2=z^2$.  It is primitive if $\gcd(x,y,z)=1$.
\end{de}
\begin{remark}
In a primitive Pythagorean triple, $\gcd(x,y)=\gcd(y,z)=\gcd(z,x)=1$, and exactly one of $x,y$ is even.
\end{remark}

\begin{thm}
The primitive Pythagorean triples with $y$ even are precisely
\[
 (x,y,z)=(r^2-s^2,\ 2rs,\ r^2+s^2),
\]
where $r>s$ are coprime positive integers of opposite parity.
\end{thm}
\begin{proof}
Let $(x,y,z)$ be primitive with $y$ even.  Then $x,z$ are odd and
\[
 \left(\frac{z+x}{2}\right)
 \left(\frac{z-x}{2}\right)
 =\left(\frac y2\right)^2.
\]
The two factors on the left are coprime, so each is a square:
$(z+x)/2=r^2$ and $(z-x)/2=s^2$.  This gives the displayed formulas.  Coprimality and opposite parity follow from primitiveness and the oddness of $x$.  Conversely, the formulas satisfy the equation, and the stated conditions make the triple primitive.
\end{proof}

\begin{remark}[rational points on the circle]
The rational points on $X^2+Y^2=1$ are parametrized by lines of rational slope through $(1,0)$:
\[
 (X,Y)=\left(\frac{m^2-1}{m^2+1},\frac{-2m}{m^2+1}\right),
 \qquad m\in\Q.
\]
Clearing denominators recovers the Pythagorean parametrization.
\end{remark}

\section{Infinite descent and local obstructions}

\begin{de}
Fermat's method of infinite descent proves nonexistence by assuming a positive integer solution, constructing a strictly smaller positive solution, and contradicting the well-ordering principle.
\end{de}

\begin{app}
\begin{enumerate}[label=(\alph*)]
\item The system
\[
 x^2+6y^2=z^2,
 \qquad
 6x^2+y^2=w^2
\]
has no positive integer solution.  Otherwise choose one minimizing $x+y+z+w$.  Adding gives
$z^2+w^2=7(x^2+y^2)$.  Since $-1$ is a nonresidue modulo $7$, this congruence forces $7\mid z,w$, and then also $7\mid x,y$.  Dividing all four entries by $7$ produces a smaller positive solution.

\item The equation
\[
 x^3+2y^3=4z^3
\]
has only the zero integer solution.  Modulo $2$, any solution has $x$ even; after substituting $x=2x_1$, one finds $y$ even, and then $z$ even.  Division by $2$ gives a smaller solution of the same equation, so infinite descent forces $x=y=z=0$.
\end{enumerate}
\end{app}

\begin{thm}[Fermat for exponent four]
The equation
\[
 x^4+y^4=z^4
\]
has no nontrivial integer solutions, where nontrivial means $xyz\ne0$.  In fact, the stronger equation
\[
 x^4+y^4=z^2
\]
has no positive integer solutions.
\end{thm}
\begin{proof}
Assume a positive solution to $x^4+y^4=z^2$ with $z$ minimal, and divide out a common factor so that it is primitive.  Exactly one of $x,y$ is even; suppose $y$ is even.  The primitive Pythagorean triple $(x^2,y^2,z)$ gives coprime $r>s$ of opposite parity with
\[
 x^2=r^2-s^2,
 \qquad
 y^2=2rs,
 \qquad
 z=r^2+s^2.
\]
The even parameter must be $s$: if $r$ were even, then $x^2=r^2-s^2\equiv3\pmod4$.  The square condition $2rs=y^2$ and coprimality now imply
\[
 r=m^2,
 \qquad
 s=2n^2.
\]
Thus
\[
 x^2=(m^2-2n^2)(m^2+2n^2).
\]
The two positive odd factors are coprime, so each is a square:
\[
 m^2-2n^2=u^2,
 \qquad
 m^2+2n^2=v^2.
\]
Then
\[
 \frac{v-u}{2}\cdot\frac{v+u}{2}=n^2.
\]
The two factors are coprime, so write them as $b^2$ and $a^2$.  Hence
$u=a^2-b^2$, $n=ab$, and
\[
 m^2=u^2+2n^2=(a^2-b^2)^2+2a^2b^2=a^4+b^4.
\]
This is another positive solution of the same form with hypotenuse $m<z$, contradicting minimality.
\end{proof}

\begin{thm}[local obstruction principle]
If an integer equation $f(x,y,z)=0$ has an integer solution, then for every modulus $m\ge2$ the congruence
\[
 f(x,y,z)\equiv0\pmod m
\]
has a solution.  Thus a single modulus with no congruence solution proves that the integer equation has no solution.
\end{thm}

\begin{app}
\begin{enumerate}[label=(\alph*)]
\item The equation
\[
 x^3+y^3-7z^3=3
\]
has no integer solution.  Modulo $7$, every cube is $0$, $1$, or $-1$, so $x^3+y^3$ can only be $-2,-1,0,1,2$, never $3$.

\item The equation
\[
 y^2=x^3+7
\]
has no integer solution.  If $x$ is even, the right side is $3$ modulo $4$.  Hence $x$ is odd.  The case $x\equiv3\pmod4$ would make the right side $2$ modulo $4$, so $x\equiv1\pmod4$.  The equation also forces $x\ge1$.  Moreover,
\[
 y^2+1=x^3+8=(x+2)(x^2-2x+4).
\]
Since $x+2\equiv3\pmod4$, some prime $p\equiv3\pmod4$ divides $x+2$ to an odd power.  Thus $-1$ is a quadratic residue modulo $p$, contradicting the first supplementary law.
\end{enumerate}
\end{app}
